\documentclass[11pt,a4paper]{article}

\usepackage[margin=2.4cm]{geometry}
\setlength{\parindent}{0pt}
\setlength{\parskip}{0.55em}
\renewcommand{\labelitemi}{--}

\begin{document}

\noindent
17 September 2026

\medskip
\noindent
The Editors\\
Special Issue on \emph{Computer Vision Systems for Document Analysis and
Recognition}\\
\emph{International Journal on Document Analysis and Recognition (IJDAR)}

\medskip
\noindent
Dear Editors,

\noindent
We are pleased to submit our manuscript, \textbf{``Reference-Conditioned
Structural Assessment of Low-Resource Chinese Calligraphy with Overlapping
Stroke Parsing and Human-Audited Spatial Scoring''}, by Xiaofan Liu, Ronghao
Zhang and Yuan Feng (Software Engineering Institute, East China Normal
University), for consideration in the above special issue. Xiaofan Liu is the
corresponding author.

\textbf{Fit to the special issue.} The special issue invites work that goes
beyond narrowly defined benchmarks and addresses document heterogeneity,
low-resource learning, and explainable evidence. Our manuscript studies a
document-analysis problem that is deliberately difficult on all three axes:
the inputs are brush-based Chinese calligraphy, the structural annotation is
scarce (769 quality-controlled samples over 40 character identities, without
reliable writer identities), and the intended output is not a single number
but inspectable evidence. We formulate
\emph{reference-conditioned structural assessment}---comparing a learner
instance against an approved reference of the same character---and report
where the approach works, where it does not, and under which conditions each
claim holds.

\textbf{Contributions.} The manuscript is organized around four research
questions.
\begin{itemize}
  \item \emph{RQ1.} We define an overlapping six-channel representation
  (five direction families that may overlap, plus a separate endpoint
  channel) and compare U-Net, DeepLabV3+ and SegFormer-B2 across three seeds
  on both a standard and a character-disjoint split, showing that the best
  in-corpus parser is not the best parser for unseen characters.
  \item \emph{RQ2.} Using controlled nuisance and structural perturbations of
  200 real references, we quantify what a bounded similarity registration
  buys: nuisance penalties fall substantially relative to no alignment while
  structural edits remain penalized. We also report a wider-search
  comparator and a channel-coverage score audit.
  \item \emph{RQ3.} Three blinded domain-qualified raters scored 150 natural
  same-character pairs. We report inter-rater reliability, the association of
  three scores, and a direct-ink control that isolates how much the scalar
  score depends on the parser. Retrospective development is explicitly
  separated from character-grouped internal validation.
  \item \emph{RQ4.} We evaluate localized diagnostic evidence under
  interventions with known ground truth and publish a failure taxonomy
  rather than an aggregate accuracy.
\end{itemize}
All reported values are computed from frozen artifacts whose SHA-256 digests
are listed in the Online Resource.

\textbf{Originality.} This submission is original work. It has not been
published previously and is not under consideration elsewhere. It does not
extend a prior conference or journal publication by the same authors, so
there is no companion paper to declare or to supply for comparison.

\textbf{Third-party material and permission.} Three figures reproduce
handwriting images from the Calli-Tongji open subset (Ouyang Xun
regular-script and Wang Xizhi running-script categories), which is
distributed under CC~BY-NC~4.0. These images are used with full attribution,
at reduced resolution, and only as illustrative examples; the dataset itself
is not redistributed, and the manuscript states this under Data Availability.
Because the licence is non-commercial while the journal is published by a
commercial publisher, we flag this explicitly here and will supply
documentation of reproduction permission if the editorial office requires
it.
%
% --- HOW TO UPDATE THIS PARAGRAPH: keep one version, delete the rest. ---
% After the dataset providers reply in writing, append:
%   We have obtained the written permission of the dataset providers to
%   reproduce these images in this article, and the correspondence is
%   uploaded as a related file.
% While a reply is still pending, append:
%   We have written to the dataset providers to request permission to
%   reproduce these images in this article and will forward their reply as
%   soon as it is received.
% If permission is refused, remove those figures and their captions and
%   re-run the figure builders before submitting.

\textbf{Format.} The manuscript is 20 pages in the IJDAR two-column format,
including references, figures and tables, and therefore within the journal's
20-page guidance. It uses a structured abstract of 228 words (Purpose,
Methods, Results, Conclusion) and six keywords. Complete editable LaTeX
source, including the class file and all figures, is uploaded as a single
archive named \texttt{OneStroke2026\_manuscript\_latex.zip}, in which
\texttt{manuscript.tex} is the main document. The Electronic Supplementary
Material (four tables and two notes) is uploaded separately as
\texttt{OneStroke2026\_ESM\_1\_submission.pdf} and is cited in the text as
Online Resource~1.

\textbf{Declarations.} All authors have approved the submission and agree to
its content. The authors declare no competing interests. The work received no
external funding. The study involved no clinical, medical, or behavioural
intervention; the handwriting corpus was commissioned without recording
writer identities, and the similarity ratings were provided under evaluator
codes by two of the authors and one collaborating specialist. Data and code
availability statements are given in the manuscript and entered in the
submission interface.

\textbf{Reviewer suggestions.} We respectfully suggest reviewers with
expertise in document image analysis and in computational calligraphy, and we
ask that experts from the authors' institution be excluded.
%
% OPTIONAL: replace the sentence above with two to four named experts and
% their institutional e-mail addresses if the editors accept suggestions.

\medskip
\noindent
We would be glad to provide any additional material the editors require.
Thank you for considering our manuscript.

\medskip
\noindent
Sincerely,\\
Xiaofan Liu\\
on behalf of all authors\\
Software Engineering Institute, East China Normal University, Shanghai, China\\
\href{mailto:10244602411@stu.ecnu.edu.cn}{10244602411@stu.ecnu.edu.cn}

\end{document}
